\section{Anhang: Reproduzierbarkeit}

Die Ausführung erfolgt mit Docker Compose; Domänenkette, Warp-Stufe und
Quality-Gates sind in Kapitel~\ref{sec:implementierung} beschrieben und gelten
für Inline-Vollauf, Matrixrunner und Replays mit frischem COLMAP gleichermaßen.
Dieser Anhang fasst nur die konstanten Parameter der PA-Versuche zusammen:

\begin{itemize}
  \item Prompt \texttt{pipe};
  \item 4096 SIFT-Merkmale, Sequential-Overlap 15, Guided Matching aus;
  \item STS 7000 Gesamtiterationen, Stage-2-Fenster 5000;
  \item Objektmasken für die Bildmetriken und feste Eval-Listen
        (210 Training / 30 Evaluation);
  \item 200.000 Mesh-Vertices, 5.000.000 Oberflächenstichproben und Seed 42;
  \item objektmaskierte PSNR-, SSIM- und LPIPS-Auswertung;
  \item Produktionsroute \texttt{original\_gs}; SuGaR-Coarse nur als
        explizite Vergleichsroute;
  \item Centerline-Modus \texttt{single}, Voxelgröße 0,1, B-Spline Grad 10,
        Eckensegmentierung deaktiviert.
\end{itemize}

\subsection*{Ausführungsmodi}

\begin{description}
  \item[Interaktiver Vollauf] führt durch Video-, COLMAP-, STS- und
    Meshroutenparameter sowie den optionalen GCP-Breakpoint.
  \item[Autopilot] fragt nach seiner Aktivierung nur Lauf-Preset und
    Objekt-Prompt ab, setzt die übrigen dokumentierten Standardwerte und
    überspringt den manuellen GCP-Breakpoint, falls keine Matrix bereitsteht.
  \item[Matrixrunner] führt deklarierte Varianten seriell aus, archiviert sie
    und bereinigt danach den Live-Workspace.
\end{description}

\subsection*{SuGaR-Fork: Versionierung}

Der mask-aware Fork liegt unter \path{third_party/SuGaR}. Parent-Gitlink,
Fork-Checkout und Docker-Parameter \texttt{SUGAR\_REF} verweisen konsistent
auf \texttt{a0fc37b} im Zweig \path{project/masked-sugar}; genau dieser
Stand wurde in allen archivierten Läufen verwendet. Die Änderungsschichten im
Einzelnen:

\begin{itemize}
  \item \texttt{48bbfdd}: maskenbewusster Semantic-Mask-Provider sowie
    maskierte RGB-/DSSIM-Supervision in den Coarse-Trainern und im Refinement,
    maskierte DN-Consistency, konfigurierbare Maskenstufen/Dilatationen,
    semantisch begrenztes UV-Baking sowie erweiterte Trainings- und
    Meshparameter;
  \item \texttt{8f4f7a2}: Poisson-Schutz bei zu kleinen Punktwolken oder
    ungültigen Normalen und Korrektur der Coarse-Iterationsprüfung für
    STS-Startzähler;
  \item \texttt{e254000}: Schalter für Gaussian-Rasterizer-Tiefe als
    kontrollierte Surface-Sampling-Diagnose;
  \item \texttt{ecda7ef}: zusätzliche Surface-Sampling-Optionen und explizite
    Background-Mesh-Steuerung;
  \item \texttt{eca4ea1}/\texttt{a0fc37b}: normalisierte feste Eval-Splits,
    Schutz gegen nicht gematchte Kameras und explizite Prüfung der
    Eval-Kameranamen.
\end{itemize}

Den vollständigen Diff ab \texttt{48bbfdd} hält die Datei
\path{sugar_fork_diff_48bbfdd_a0fc37b.diff} im digitalen Anlagenband fest.
Repository-Commit, Submodul-Commit, Docker-Image-Version und Hardware werden
zusätzlich im finalen Laufmanifest dokumentiert.
